\documentclass[11pt]{article}
\usepackage[margin=1in]{geometry}
\usepackage{amsmath,amssymb,booktabs,array,longtable}
\usepackage{graphicx}
\graphicspath{{./}{report_figures/}}

\usepackage{mathptmx} % Times New Roman for text and math
\usepackage{hyperref}
\usepackage{enumitem}
\usepackage{titlesec}
\usepackage[table]{xcolor}
\usepackage{float}
\usepackage{placeins}
\usepackage{caption}

\captionsetup{font=small,labelfont=bf,justification=centering}

\hypersetup{
    colorlinks=true,
    linkcolor=blue,
    urlcolor=blue,
    citecolor=blue
}

\setlength{\parskip}{4pt}
\setlength{\parindent}{0pt}

\begin{document}

\begin{center}
\Large\textbf{Shiv Nadar University Chennai}\\[2mm]
\end{center}

\renewcommand{\arraystretch}{1.4}

\begin{tabular}{|p{3.5cm}|p{8cm}|p{2cm}|}
\hline
\rowcolor{blue!8}
Name of Student & \textbf{KR Srinidhi Rajamane} & Reg.\ No.\\[1mm]
\rowcolor{blue!8}
 & & \textbf{24011101059}\\
\hline
Degree \& Branch &
\begin{tabular}[t]{@{}l@{}}
B.Tech Artificial Intelligence \& Data Science\\
Section A
\end{tabular}
&
Semester V\\
\hline
Subject Code \& Name & CS3807 -- Deep Learning Laboratory & AY: 2026--27\\ \hline
\end{tabular}

\vspace{3mm}

\begin{center}
\Large\textbf{Experiment 1}\\
\large\textbf{Implementation of a Single Layer Perceptron for Binary Classification}
\end{center}

%----------------------------------------------------------

\section*{1. Objective}

The objective of this experiment is to understand the concept of an artificial neuron and implement a Single Layer Perceptron from scratch for binary classification. Students will learn the perceptron learning algorithm, understand the importance of activation functions, visualize the learning process, and evaluate the classifier using a real-world dataset.

%----------------------------------------------------------

\section*{2. Tasks to be Performed}

\begin{enumerate}[leftmargin=*]
\item Download and understand the dataset.
\item Perform Exploratory Data Analysis (EDA).
\item Preprocess the dataset.
\item Implement a Single Layer Perceptron from scratch.
\item Train the model using the perceptron learning rule.
\item Evaluate the classifier using standard performance metrics.
\item Analyze the effect of different learning rates.
\item Compare the implementation with Scikit-learn's Perceptron (optional).
\end{enumerate}

%----------------------------------------------------------

\section*{3. Background Theory}

An Artificial Neuron is the fundamental building block of a neural network. It receives one or more input features, multiplies them by corresponding weights, adds a bias term, and applies an activation function to determine the output.

The Single Layer Perceptron is the earliest supervised learning algorithm capable of solving linearly separable binary classification problems.

\subsection*{Mathematical Formulation}

Weighted Sum

\[
z=\mathbf{w}^{T}\mathbf{x}+b
\]

where

\begin{itemize}
\item $\mathbf{x}$ : Input vector
\item $\mathbf{w}$ : Weight vector
\item $b$ : Bias
\item $z$ : Linear combination
\end{itemize}

Prediction

\[
\hat y=f(z)
\]

Weight Update Rule

\[
\mathbf{w}_{new}
=
\mathbf{w}_{old}
+
\eta
(y-\hat y)
\mathbf{x}
\]

Bias Update

\[
b_{new}
=
b_{old}
+
\eta
(y-\hat y)
\]

where $\eta$ denotes the learning rate.

%----------------------------------------------------------

\section*{4. Activation Functions}

An activation function determines the output of a neuron based on the weighted sum of its inputs. It introduces non-linearity, enabling neural networks to learn complex patterns. Although the perceptron uses only the Step Activation Function, modern deep learning models employ differentiable activation functions for efficient optimization through backpropagation.

\subsection*{Step Activation Function}

\[
f(z)=
\begin{cases}
1,&z\ge0\\
0,&z<0
\end{cases}
\]

The Step Activation Function produces a binary output and is therefore suitable only for linearly separable binary classification problems.

\begin{center}
\begin{tabular}{|p{2.5cm}|p{3cm}|p{6cm}|}
\hline
\textbf{Activation} & \textbf{Output Range} & \textbf{Typical Applications}\\
\hline
Step & $\{0,1\}$ & Classical Perceptron\\
\hline
Sigmoid & $(0,1)$ & Binary Classification Output Layer\\
\hline
Tanh & $(-1,1)$ & Hidden Layers (Traditional Networks)\\
\hline
ReLU & $[0,\infty)$ & Hidden Layers in CNNs and MLPs\\
\hline
Leaky ReLU & $(-\infty,\infty)$ & Avoiding Dying ReLU\\
\hline
Softmax & $(0,1)$ & Multi-class Classification Output Layer\\
\hline
\end{tabular}
\end{center}

\textbf{Note:} This experiment uses only the Step Activation Function. The remaining activation functions will be studied in subsequent experiments involving Multilayer Perceptrons and Deep Neural Networks.

%----------------------------------------------------------

\section*{5. Dataset Description}

\textbf{Dataset:} Banknote Authentication Dataset

\textbf{Source:} UCI Machine Learning Repository

\textbf{Download Link:}

\url{https://archive.ics.uci.edu/dataset/267/banknote+authentication}

\begin{center}
\begin{tabular}{|l|l|}
\hline
Instances & 1372\\
\hline
Features & 4 Numerical Features\\
\hline
Classes & 2\\
\hline
Missing Values & None\\
\hline
Task & Binary Classification\\
\hline
\end{tabular}
\end{center}

\textbf{Feature Description}

\begin{itemize}
\item Variance
\item Skewness
\item Curtosis
\item Entropy
\end{itemize}

Target Class

\begin{itemize}
\item 0 -- Authentic Banknote
\item 1 -- Forged Banknote
\end{itemize}

%----------------------------------------------------------

\section*{6. Experimental Procedure}

\textbf{Task 1: Dataset Exploration}

\begin{itemize}
\item Load the dataset.
\item Display the first five samples.
\item Determine dataset dimensions.
\item Identify missing values.
\item Display descriptive statistics.
\end{itemize}

\textbf{Expected Output}

Dataset summary and statistical information.

\vspace{2mm}

\textbf{Task 2: Exploratory Data Analysis}

Generate

\begin{itemize}
\item Histograms
\item Correlation Heatmap
\item Scatter Plot
\item Boxplots
\end{itemize}

\vspace{2mm}

\textbf{Task 3: Data Preprocessing}

\begin{itemize}
\item Normalize all numerical features.
\item Split the dataset into Training (80\%) and Testing (20\%).
\end{itemize}

\vspace{2mm}

\textbf{Task 4: Perceptron Implementation}

Implement from scratch

\begin{itemize}
\item Weight Initialization
\item Bias Initialization
\item Step Activation Function
\item Forward Propagation
\item Perceptron Learning Rule
\end{itemize}

\vspace{2mm}

\textbf{Task 5: Model Training}

Display

\begin{itemize}
\item Epoch Number
\item Misclassified Samples
\item Updated Weights
\item Updated Bias
\end{itemize}

\vspace{2mm}

\textbf{Task 6: Model Evaluation}

Compute

\begin{itemize}
\item Accuracy
\item Precision
\item Recall
\item F1-score
\item Confusion Matrix
\end{itemize}

%----------------------------------------------------------

\section*{7. Mandatory Plots}

\begin{longtable}{|p{3.4cm}|p{3.4cm}|p{6cm}|}
\hline
\textbf{Plot} & \textbf{Purpose} & \textbf{Expected Inference}\\
\hline
Feature Histograms & Distribution Analysis & Identify skewness and outliers\\
\hline
Correlation Heatmap & Feature Relationship & Observe highly correlated features\\
\hline
Scatter Plot & Class Distribution & Determine whether classes are approximately linearly separable\\
\hline
Training Error vs Epoch & Learning Behaviour & Verify convergence of the perceptron\\
\hline
Weight Evolution & Parameter Analysis & Observe stabilization of learned weights\\
\hline
Bias Evolution & Parameter Analysis & Understand the role of the bias term\\
\hline
Confusion Matrix & Performance Evaluation & Interpret TP, FP, TN and FN\\
\hline
Learning Rate Comparison & Hyperparameter Study & Compare convergence behaviour for different learning rates\\
\hline
Decision Boundary (Optional) & Visualization & Illustrate the separating hyperplane using two selected features\\
\hline
\end{longtable}

%----------------------------------------------------------

\section*{8. Performance Tables}

\textbf{Training Summary}

\begin{tabular}{|p{6cm}|p{7cm}|}
\hline
Dataset Size & 1372 samples\\
\hline
Train/Test Split & 80\% / 20\% (1097 train / 275 test)\\
\hline
Learning Rate & 0.01\\
\hline
Epochs & 50 (ran for the full budget; did not reach zero error)\\
\hline
Final Weights & Variance: $-0.1637$, Skewness: $-0.1767$, Curtosis: $-0.1895$, Entropy: $0.0040$\\
\hline
Final Bias & $0.2400$\\
\hline
Accuracy & $0.9709$ (97.09\%)\\
\hline
Precision & $0.9385$ (93.85\%)\\
\hline
Recall & $1.0000$ (100.00\%)\\
\hline
F1-score & $0.9683$ (96.83\%)\\
\hline
\end{tabular}

\vspace{3mm}

\textbf{Epoch-wise Learning} \\[1mm]
\textit{(Weight 1 = Variance, Weight 2 = Skewness; selected epochs shown for brevity)}

\begin{tabular}{|c|c|c|c|c|}
\hline
Epoch & Errors & Weight 1 & Weight 2 & Bias\\
\hline
1  & 175 & $-0.0744$ & $-0.0491$ & $0.0700$\\
5  & 32  & $-0.0964$ & $-0.0952$ & $0.1200$\\
10 & 36  & $-0.1125$ & $-0.1215$ & $0.1500$\\
20 & 20  & $-0.1231$ & $-0.1465$ & $0.1800$\\
30 & 22  & $-0.1389$ & $-0.1609$ & $0.2000$\\
40 & 18  & $-0.1499$ & $-0.1795$ & $0.2200$\\
47 & 16  & $-0.1642$ & $-0.1847$ & $0.2300$\\
50 & 27  & $-0.1637$ & $-0.1767$ & $0.2400$\\
\hline
\end{tabular}

%----------------------------------------------------------

\section*{9. Additional Tasks}

\subsection*{9.1 Step vs. Sigmoid Activation}

\begin{tabular}{|p{4cm}|p{5.5cm}|p{5.5cm}|}
\hline
\textbf{Property} & \textbf{Step} & \textbf{Sigmoid}\\
\hline
Output & Hard $\{0,1\}$ & Smooth $(0,1)$\\
\hline
Derivative & Zero everywhere, undefined at $z=0$ & Well-defined, non-zero almost everywhere\\
\hline
Gradient flow & None -- blocks backpropagation & Enables gradient-based learning\\
\hline
Confidence info & None (hard decision only) & Yes, output is a probability-like score\\
\hline
\end{tabular}

The Step function is not used in modern deep learning because it is non-differentiable, so gradient descent and backpropagation cannot flow a useful error signal back through it. A single perceptron only ever needs a hard yes/no decision, but a multi-layer network needs every layer's error to be propagated backward using the chain rule, which requires a smooth, differentiable activation such as Sigmoid, Tanh or ReLU. 

\vspace{3mm}
\subsection*{9.2 Comparison with Scikit-learn's Perceptron}

The from scratch implementation matched the Scikit-learn's accuracy, precision, recall and F1 score.

\vspace{3mm}
\subsection*{9.3 Effect of Learning Rate}

This was repeated for various learning rates, each trained for 50 epochs; the resulting convergence curves are plotted and the final error counts are tabulated. In short the larger step sizes overshoot the decision boundary on data that is not perfectly linearly separable instead of settling into it.

\vspace{3mm}
\subsection*{9.4 Why a Single Layer Perceptron Cannot Solve XOR}

XOR is not linearly separable. Since a single-layer perceptron can only represent one linear decision boundary, it is mathematically incapable of solving XOR. Solving XOR requires at least one hidden layer with a non-linear activation function which remaps the inputs into a new feature space in which the classes become linearly separable.

\vspace{3mm}
\subsection*{9.5 Effect of Feature Normalization on Convergence}

Before scaling Skewness has a much larger numeric range than Entropy, so the weighted sum would be dominated by whichever feature happens to have the largest raw magnitude, not by which feature is actually most informative for classification. This slows and destabilizes convergence.

%----------------------------------------------------------

\section*{10. Source Code}

The complete implementation was carried out in Python using NumPy, Pandas, Matplotlib, Seaborn and Scikit-learn. The full, executable source code (organized task-wise, exactly as it was run) is available on GitHub:

\begin{center}
\href{https://github.com/Srinidhi-Krishna/Deep-Learning-Lab/blob/main/Lab_1/DL_LAB_1.ipynb}{\textbf{https://github.com/Srinidhi-Krishna/Deep-Learning-Lab/blob/main/Lab\_1/DL\_LAB\_1.ipynb}}
\end{center}

%----------------------------------------------------------

\section*{11. Experimental Results}

\subsection*{11.1 Dataset Exploration (Task 1)}

The loaded data set is fine. It has 1372 rows and 5 columns (4 features + class label). Skewness has the widest spread, while Entropy is the most tightly packed. 

\subsection*{11.2 Data Preprocessing (Task 3)}

After \texttt{MinMaxScaler}, all four features got squeezed into $[0,1]$. 

\subsection*{11.3 Model Training (Task 5)}

Trained for a fixed 50 epochs. Though initial epochs started with considerable errors, it dropped fast and then just kept bouncing between for the rest of training instead of settling down.  

\subsection*{11.4 Model Evaluation (Task 6)}

Final-epoch weights on the test set gave:

\begin{center}
\begin{tabular}{|l|c|}
\hline
\textbf{Metric} & \textbf{Value}\\
\hline
Accuracy  & 0.9709 (97.09\%)\\
Precision & 0.9385 (93.85\%)\\
Recall    & 1.0000 (100.00\%)\\
F1-score  & 0.9683 (96.83\%)\\
\hline
\end{tabular}
\end{center}
\vspace{2cm}
Confusion matrix on the 275 test samples:
\vspace{0.2cm}
\[
\begin{bmatrix} 145 & 8 \\ 0 & 122 \end{bmatrix}
\]

145 authentic notes were correctly caught, 8 authentic notes were wrongly flagged as forged and zero forged notes slipped through as authentic. Hence the recall on the forged class is a perfect 1.0. 

\vspace{0.5cm}
\subsection*{11.5 Effect of Learning Rate (Task 7)}

Tried three rates, each trained fresh for 50 epochs:
\vspace{0.5cm}
\begin{center}
\begin{tabular}{|c|c|c|}
\hline
\textbf{Learning Rate ($\eta$)} & \textbf{Epochs Run} & \textbf{Final Misclassified}\\
\hline
0.001 & 50 & 16\\
0.01  & 50 & 27\\
0.1   & 50 & 27\\
\hline
\end{tabular}
\end{center}

I expected the smallest rate to be the slowest one, but surprisingly it ended up with the fewest errors. Small steps mean that the weight vector moves cautiously and does not overshoot.

\vspace{0.5cm}
\subsection*{11.6 Comparison with Scikit-learn (Task 9, Optional)}
\vspace{0.5cm}
Ran scikit-learn's \texttt{Perceptron} with the same settings on the same split and got identical numbers to the from-scratch version for accuracy, precision, recall and F1 score.

%----------------------------------------------------------

\vspace{2mm}
\section*{12. Generated Plots with Interpretation}

\begin{figure}[H]
\centering
\includegraphics[width=0.62\textwidth]{Fig1_Histograms}
\caption{Frequency distribution histograms of Variance, Skewness, Curtosis and Entropy.}
\end{figure}

\noindent\textbf{Interpretation:} None of the four features appears bell-shaped. The raw features are not normally distributed. This is not a problem for a perceptron, but this is why I decided to go with min-max scaling instead of standardization.

\vspace{4mm}

\begin{figure}[H]
\centering
\includegraphics[width=0.55\textwidth]{Fig2_CorrelationHeatmap}
\caption{Pearson correlation heatmap among the four features and the class label.}
\end{figure}

\noindent\textbf{Interpretation:} Of the four, Variance has the strongest correlation to Class, which is consistent with the final weights. 

\FloatBarrier

\begin{figure}[H]
\centering
\includegraphics[width=0.55\textwidth]{Fig3_ScatterPlot}
\caption{Scatter plot of Variance versus Skewness, colour-coded by class.}
\end{figure}

\noindent\textbf{Interpretation:} This plot is what convinced me that a linear classifier could actually work here as authentic and forged points sit mostly in different regions.

\vspace{4mm}

\begin{figure}[H]
\centering
\includegraphics[width=0.62\textwidth]{Fig4_Boxplots}
\caption{Boxplots of each feature grouped by class.}
\end{figure}

\noindent\textbf{Interpretation:} Same story as the scatter plot with median Variance and Skewness for authentic notes sit above forged notes. 

\FloatBarrier

\begin{figure}[H]
\centering
\includegraphics[width=0.58\textwidth]{Fig5_ConvergenceCurve}
\caption{Misclassified training samples per epoch (learning rate 0.01).}
\end{figure}

\noindent\textbf{Interpretation:} Errors drop fast at first and then flatten into a noisy plateau for the rest of the training. It never hits zero, which means that the four features together aren't perfectly linearly separable.

\vspace{4mm}

\begin{figure}[H]
\centering
\includegraphics[width=0.58\textwidth]{Fig7_WeightEvolution}
\caption{Evolution of the four feature weights across training epochs.}
\end{figure}

\noindent\textbf{Interpretation:} All four weights move fast early on and then level off into small oscillations instead of a flat line, matching the fact that error never hits zero. 

\FloatBarrier

\begin{figure}[H]
\centering
\includegraphics[width=0.55\textwidth]{Fig8_BiasEvolution}
\caption{Evolution of the bias term across training epochs.}
\end{figure}

\noindent\textbf{Interpretation:} The bias climbs in a rough staircase pattern from 0 up to about 0.24 and it's still climbing at epoch 50, which means the model is still making corrections right up to the end.

\vspace{4mm}

\begin{figure}[H]
\centering
\includegraphics[width=0.5\textwidth]{Fig6_ConfusionMatrix}
\caption{Confusion matrix on the held-out test set (275 samples).}
\end{figure}

\noindent\textbf{Interpretation:} Only 8 out of 275 test notes were misclassified, and all 8 were authentic notes flagged as forged notes but none of the forged notes got through as authentic. For currency verification that is a good error pattern to have, since a false alarm on a genuine note is a lot less costly than letting a fake one pass.

\FloatBarrier

\begin{figure}[H]
\centering
\includegraphics[width=0.58\textwidth]{Fig9_LearningRateComparison}
\caption{Convergence behaviour of the perceptron under three different learning rates.}
\end{figure}

\noindent\textbf{Interpretation:} All three curves drop rapidly at first but settle at different noise floors. Smaller steps here land closer to a good boundary, larger steps keep the weight vector jumping around more.

\vspace{4mm}

\begin{figure}[H]
\centering
\includegraphics[width=0.55\textwidth]{Fig10_DecisionBoundary}
\caption{Decision boundary learned by a 2-feature perceptron (Variance vs Skewness).}
\end{figure}

\noindent\textbf{Interpretation:} With just two features, the boundary is seen having the most authentic points on one side, most forged on the other, with a few points from each class right on or just past the line. 

\FloatBarrier
%----------------------------------------------------------

\section*{13. Discussion}

The from-scratch perceptron does well here with a test accuracy of  97.09\% and 100\% recall on the forged class. A fake note never passed through as genuine on the test split.

The learning-rate points out that the smallest rate ($\eta=0.001$) ended up with fewer errors than the larger ones, which goes against the usual "bigger step = faster learning" intuition. 

Normalization mattered too as Skewness has a much bigger raw scale than Entropy, so without scaling, Skewness would dominate the weighted sum.

%----------------------------------------------------------

\section*{14. Conclusion}

A Single Layer Perceptron is implemented from scratch and it is tested on the UCI Banknote Authentication dataset. This resulted in  97.09\% accuracy with 93.85\% precision, 100\% recall and 96.83\% F1 on the test set, and the output matched scikit-learn's built-in \texttt{Perceptron} exactly under the same hyperparameters.

The training curves showed that the four features are only approximately linearly separable, so the model settles into a noisy low-error region instead of reaching zero, and the learning-rate experiment showed that a smaller step size can actually converge better than a larger one on this kind of data.


%----------------------------------------------------------

\section*{15. References}

\begin{enumerate}
\item F. Rosenblatt, ``The Perceptron,'' \emph{Psychological Review}, 1958.
\item I. Goodfellow, Y. Bengio and A. Courville, \emph{Deep Learning}, MIT Press, 2016.
\item C. M. Bishop, \emph{Pattern Recognition and Machine Learning}, Springer, 2006.
\item S. Haykin, \emph{Neural Networks and Learning Machines}, Pearson, 2009.
\item UCI Machine Learning Repository -- Banknote Authentication Dataset.
\item Scikit-learn Documentation: Perceptron.
\end{enumerate}

\end{document}
